
\subsubsection{Quantitative Findings}

The simulation produced the following measurements:

\begin{table}[htbp]
\centering
\resizebox{\columnwidth}{!}{%
\begin{tabular}{ll}
\toprule
Metric & Value \\
\midrule
ECE before calibration & 0.5267 \\
ECE after calibration & 0.0798 \\
ECE reduction & 84.8\% \\
Absolute ECE decrease & 0.4469 \\
Optimal temperature T* & 2512.712 \\
Pass@1 (calibration) & 36.00\% (unchanged) \\
Pass@1 (validation) & 42.05\% (unchanged) \\
\bottomrule
\end{tabular}
}
\end{table}

\textbf{Interpretation:} The simulation confirms that (1) temperature scaling optimization converges, (2) ECE decreases after calibration, (3) ranking-based accuracy is preserved. The specific ECE values and temperature magnitude are simulation artifacts and should not be interpreted as properties of real code generation models.

\subsubsection{Optimization Convergence}

LBFGS optimization converged monotonically over 200 iterations with no oscillation or divergence. Negative log-likelihood decreased from 0.8234 to 0.2156. This validates that the optimization procedure is stable for single-parameter temperature scaling.

\textbf{Temperature magnitude artifact:} The optimal temperature T* = 2512.71 is substantially higher than values reported for image classification (typically 0.5-3.0). This is an artifact of binary logit representation. Real Code Llama experiments with ~32K-dimensional logits would yield different temperature values. The key validation is that optimization converges, not the specific converged value.

\subsubsection{Reliability Diagrams}

The reliability diagram shows predicted confidence (x-axis) versus empirical accuracy (y-axis). Before calibration, predictions deviate from the diagonal (perfect calibration line). After calibration, predictions align closer to the diagonal.

This visualization confirms that the plotting code functions correctly and that ECE reduction corresponds to improved diagonal alignment. The specific degree of alignment in simulation should not be interpreted as representative of real model behavior.

\subsubsection{Per-Bin Analysis}

Calibration error by confidence bin:

\begin{table}[htbp]
\centering
\resizebox{\columnwidth}{!}{%
\begin{tabular}{llll}
\toprule
Bin Range & Before & After & Reduction \\
\midrule
0.0-0.1 & 0.023 & 0.009 & 62.0\% \\
0.1-0.2 & 0.046 & 0.012 & 73.0\% \\
0.2-0.3 & 0.068 & 0.015 & 78.6\% \\
0.3-0.4 & 0.089 & 0.017 & 81.3\% \\
0.4-0.5 & 0.123 & 0.020 & 84.0\% \\
0.5-0.6 & 0.157 & 0.023 & 85.1\% \\
0.6-0.7 & 0.189 & 0.029 & 84.7\% \\
0.7-0.8 & 0.235 & 0.035 & 85.3\% \\
0.8-0.9 & 0.346 & 0.046 & 86.8\% \\
0.9-1.0 & 0.568 & 0.068 & 88.1\% \\
\bottomrule
\end{tabular}
}
\end{table}

The simulation shows larger error reductions in high-confidence bins. This is expected from the data generation procedure (high-confidence predictions were assigned lower accuracy to create baseline miscalibration). The pattern validates that per-bin error analysis code executes correctly.

\subsubsection{Summary of Validated Technical Components}

The simulation successfully validated:
\begin{enumerate}
\item ✓ LBFGS optimization converges for temperature parameter
\item ✓ ECE computation produces stable measurements
\item ✓ Calibration/validation data splitting is implementable
\item ✓ Pass@1 accuracy is preserved (order-preserving property)
\item ✓ Reliability diagrams are generated correctly
\item ✓ Per-bin error analysis executes without errors
\end{enumerate}

The simulation does not validate:
\begin{itemize}
\item Actual calibration quality of Code Llama 7B
\item Generalization to real code generation distributions
\item Optimal temperature value for real models
\end{itemize}
